\chapter{Schéma de la base de données}

\section{Modèle Entité-Relation}

Cette annexe présente les détails complets du schéma de la base de données HayaBook tel que défini dans le fichier \texttt{prisma/schema.prisma}. Ce schéma est utilisé par l'ORM Prisma pour générer automatiquement le client TypeScript typé permettant d'interagir avec la base de données PostgreSQL.

\subsection{Modèle \texttt{User}}

Le modèle \texttt{User} est la table centrale du système. Il représente tous les acteurs de la plateforme (clients, prestataires et administrateurs).

\begin{lstlisting}[language=SQL, caption={Champs principaux du modèle User (Prisma Schema)}]
model User {
  id                String           @id @default(uuid())
  email             String?          @unique
  phone             String?          @unique
  passwordHash      String
  firstName         String?
  lastName          String?
  bio               String?
  role              Role             @default(CLIENT)
  profileImage      String?
  isActive          Boolean          @default(true)
  isVerified        Boolean          @default(false)
  isProfileComplete Boolean          @default(false)
  isSuspended       Boolean          @default(false)
  suspensionReason  String?
  deletedAt         DateTime?
  createdAt         DateTime         @default(now())
  updatedAt         DateTime         @updatedAt
}
\end{lstlisting}

\subsection{Modèle \texttt{Booking}}

Le modèle \texttt{Booking} enregistre chaque réservation avec ses références aux entités liées et son statut dans le cycle de vie.

\begin{lstlisting}[language=SQL, caption={Champs principaux du modèle Booking (Prisma Schema)}]
model Booking {
  id                String        @id @default(uuid())
  clientId          String
  providerProfileId String
  serviceId         String
  serviceOptionId   String?
  date              DateTime
  startTime         String
  endTime           String
  durationMinutes   Int
  status            BookingStatus @default(CONFIRMED)
  price             Float
  notes             String?
  createdAt         DateTime      @default(now())
  updatedAt         DateTime      @updatedAt
}

enum BookingStatus {
  PENDING
  CONFIRMED
  IN_PROGRESS
  COMPLETED
  CANCELLED_BY_CLIENT
  CANCELLED_BY_PROVIDER
  NO_SHOW
}
\end{lstlisting}

\subsection{Relations entre les modèles}

Les principales relations entre les modèles sont :

\begin{itemize}
    \item \texttt{User} $\rightarrow$ \texttt{ProviderProfile} : Relation 1:1. Un utilisateur de rôle \texttt{PROVIDER} possède un seul profil prestataire.
    \item \texttt{ProviderProfile} $\rightarrow$ \texttt{Service} : Relation 1:N. Un prestataire peut proposer plusieurs services.
    \item \texttt{Service} $\rightarrow$ \texttt{ServiceOption} : Relation 1:N. Un service peut avoir plusieurs variantes tarifaires.
    \item \texttt{User} $\rightarrow$ \texttt{Booking} (clientBookings) : Relation 1:N. Un client peut avoir plusieurs réservations.
    \item \texttt{ProviderProfile} $\rightarrow$ \texttt{Booking} : Relation 1:N. Un prestataire peut recevoir plusieurs réservations.
    \item \texttt{Booking} $\rightarrow$ \texttt{Review} : Relation 1:1. Une réservation complétée peut donner lieu à un seul avis.
\end{itemize}

\section{Endpoints de l'API REST}

Le tableau suivant liste les principaux endpoints exposés par le serveur HayaBook.

\begin{table}[h]
\centering
\caption{Endpoints principaux de l'API REST HayaBook.}
\label{tab:endpoints}
\begin{tabular}{|l|l|l|l|}
\hline
\textbf{Méthode} & \textbf{Endpoint} & \textbf{Accès} & \textbf{Description} \\
\hline
POST & \texttt{/api/auth/register} & Public & Inscription d'un utilisateur \\
\hline
POST & \texttt{/api/auth/login} & Public & Connexion \\
\hline
POST & \texttt{/api/auth/verify-otp} & Public & Vérification OTP \\
\hline
PATCH & \texttt{/api/auth/profile} & Auth & Mise à jour du profil \\
\hline
GET & \texttt{/api/providers/all} & Public & Liste des prestataires \\
\hline
POST & \texttt{/api/bookings} & Auth & Créer une réservation \\
\hline
GET & \texttt{/api/bookings/client} & Auth & Réservations du client \\
\hline
PATCH & \texttt{/api/bookings/:id/status} & Auth & Mise à jour de statut \\
\hline
POST & \texttt{/api/messages} & Auth & Envoyer un message \\
\hline
GET & \texttt{/api/search} & Public & Recherche globale \\
\hline
GET & \texttt{/api/admin/stats} & Admin & Statistiques plateforme \\
\hline
DELETE & \texttt{/api/admin/users/:id} & Admin & Suppression douce \\
\hline
\end{tabular}
\end{table}